\slcol{अर्जुन उवाच —\\
\Index{संन्यासं कर्मणां} कृष्ण पुनर्योगं च शंससि ।\\
यच्छ्रेय एतयोरेकं तन्मे ब्रूहि सुनिश्चितम् ॥ ५.१ ॥}
\translate{Arjuna uvāca —\\
saṁnyāsaṁ karmaṇāṁ kṛṣṇa punaryōgaṁ ca śaṁsasi ।\\
yacchrēya ētayōrēkaṁ tanmē brūhi suniścitam ॥ 5.1 ॥}
\cquote{Arjuna said:
O Krishna, You praise both \textbf{Sannyasa} (renunciation of actions) and the path of \textbf{Karma Yoga} (selfless action). Kindly clarify which of these two is better.}


\slcol{श्रीभगवानुवाच —\\
\Index{संन्यासः कर्मयोगश्च} निःश्रेयसकरावुभौ ।\\
तयोस्तु कर्मसंन्यासात्कर्मयोगो विशिष्यते ॥ ५.२ ॥}
\translate{śrībhagavānuvāca —\\
saṁnyāsaḥ karmayōgaśca niḥśrēyasakarāvubhau ।\\
tayōstu karmasaṁnyāsātkarmayōgō viśiṣyatē ॥ 5.2 ॥}
\cquote{Sri Bhagavan said:\\
Both Sannyasa* and Karma Yoga* lead to highest liberation. However, of these two, Karma Yoga* is superior to renunciation.}

\newpage
\begin{mananam}{\mananamfont {Mananam sloka - 2}} 
\mananamtext Do I take on more activities than is necessary? Am I able to withdraw from activities that are not useful or necessary? On the other hand, do I tend to relinquish my necessary duties and responsibilities, especially when under stress or faced with obstacles? How can I learn to balance what I need to do and what I don’t have to? How can I engage fully in actions while maintaining inner detachment?  
\end{mananam}
\WritingHand\enspace{\selfReflect\small{Self Reflection}}
\begin{inspiration}{\mananamfont Inspiration}
\mananamtext Of the two spiritual attitudes towards activity—renunciation and actions without expectations of rewards—the latter offers a greater benefit to the majority who have responsibilities and duties which cannot be given up. However, the wiser way is to know when to push oneself towards action and when to withdraw from activities that are not necessary or helpful on one’s spiritual path. Everyone should cultivate this wisdom, knowing when to swing into action and when to withdraw from it.
\end{inspiration}
\newpage


\slcol{\Index{ज्ञेयः स नित्यसंन्यासी} यो न द्वेष्टि न काङ्क्षति ।\\
निर्द्वन्द्वो हि महाबाहो सुखं बन्धात्प्रमुच्यते ॥ ५.३ ॥}
\translate{jñēyaḥ sa nityasaṁnyāsī yō na dvēṣṭi na kāṅkṣati ।\\
nirdvandvō hi mahābāhō sukhaṁ bandhātpramucyatē ॥ 5.3 ॥}
\cquote{O \textbf{Mahābāhō} (Arjuna, mighty-armed), a person who neither harbors hatred nor desires and remains free from dualities is to be regarded as the eternal renunciate. Such a person is easily liberated from the bondage of material existence.}

\slcol{\Index{साङ्ख्ययोगौ पृथग्बालाः} प्रवदन्ति न पण्डिताः ।\\
एकमप्यास्थितः सम्यगुभयोर्विन्दते फलम् ॥ ५.४ ॥}
\translate{sāṅ‍khyayōgau pr̥thagbālāḥ pravadanti na paṇḍitāḥ ।\\
ēkamapyāsthitaḥ samyagubhayōrvindatē phalam ॥ 5.4 ॥}
\cquote{The ignorant perceive the \textbf{Sāṅkhya Yoga} (Jñāna Yoga, the paths of knowledge) and Karma Yoga* (the path of selfless action) as separate and distinct. However, the wise understand that being firmly established in either path leads to the same ultimate result.}


\slcol{\Index{यत्साङ्ख्यैप्राप्यते स्थानं} तद्योगैरपि गम्यते ।\\
एकं साङ्ख्यं च योगं च यः पश्यति स पश्यति ॥ ५.५ ॥}
\translate{yatsāṅ‍khyaiḥ prāpyatē sthānaṁ tadyōgairapi gamyatē ।\\
ēkaṁ sāṅ‍khyaṁ ca yōgaṁ ca yaḥ paśyati sa paśyati ॥ 5.5 ॥}
\cquote{The state achieved through \textbf{Sāṅkhya Yoga} can also be attained through \textbf{Karma Yoga}. A person who recognizes that \textbf{Sāṅkhya Yoga} and Karma Yoga both lead to the same goal has truly grasped their essence.}

\slcol{\Index{संन्यासस्तु महाबाहो} दुःखमाप्तुमयोगतः ।\\
योगयुक्तो मुनिर्ब्रह्म नचिरेणाधिगच्छति ॥ ५.६ ॥}
\translate{saṁnyāsastu mahābāhō duḥkhamāptumayōgataḥ ।\\
yōgayuktō munirbrahma nacirēṇādhigacchati ॥ 5.6 ॥}
\cquote{O \textbf{Mahābāhō} (Arjuna, mighty-armed), the path of \textbf{Sannyasa} (renunciation) is difficult to attain without the practice of \textbf{Karma Yoga} (the path of selfless action). However, a sage who is dedicated to Karma Yoga quickly attains the Supreme Brahman.}

\newpage
\begin{mananam}{\mananamfont {Mananam sloka - 3}} 
\mananamtext Am I aware of my likes and dislikes and how they condition my choices and decisions in life? Am I able to continue with dutiful actions even when I harbour unconscious dislikes, and discontinue unproductive actions when I harbour a strong like or attachment towards them? How can I learn to function with interest but without attachment?
\end{mananam}
\WritingHand\enspace{\selfReflect\small{Self Reflection}}
\begin{inspiration}{\mananamfont Inspiration}
\mananamtext Inner renunciation is an essential life skill, regardless of whether one is a monk or a householder. Such a one is an eternal sannyasi and has perfected the sense of detachment in the mind. He or she may have preferences of a certain lifestyle and habits that are conducive for sadhana, but inwardly they harbour neither likes nor dislikes towards anything or anyone. When circumstances warrant it, they are quick to drop anything that they have taken on, and move on with life.
\end{inspiration}
\newpage

\newpage
\begin{mananam}{\mananamfont {Mananam sloka - 6}} 
\mananamtext Do I wish to pursue the spiritual life whole heartedly? Am I ready to give up all activities at least most of them in pursuit of spirituality? Am I qualified to renounce worldly activities or to take up retirement whenever it is due, or even at an earlier juncture? Do I have a spiritual discipline or sadhana that I can commit to? How do I see myself engaging my time productively and toward the spiritual direction in which I wish to proceed?
\end{mananam}
\WritingHand\enspace{\selfReflect\small{Self Reflection}}
\begin{inspiration}{\mananamfont Inspiration}
\mananamtext To pursue the spiritual path wholeheartedly, one must commit to a sadhana. To renounce worldly life, partly or wholly, without having a disciplined life is a sure cause of sensory indulgence and laziness (otherwise called the tamasic life). A traditional spiritual path, comprising of seva, worship, meditation and study is the best bet for many who wish to live the spiritual life wholly, regardless of the stage of their life.
\end{inspiration}
\newpage

\slcol{\Index{योगयुक्तो विशुद्धात्मा} विजितात्मा जितेन्द्रियः ।\\
सर्वभूतात्मभूतात्मा कुर्वन्नपि न लिप्यते ॥५. ७ ॥}
\translate{yōgayuktō viśuddhātmā vijitātmā jitēndriyaḥ ।\\
sarvabhūtātmabhūtātmā kurvannapi na lipyatē ॥5. 7 ॥}
\cquote{A person who is united with Karma Yoga, whose soul is purified, who has gained mastery over the Self, who has conquered the senses, and who realises his Self as the Self in all beings,remains untouched by actions, even while performing them.}

\slcol{\Index{नैव किञ्चित्करोमीति} युक्तो मन्येत तत्त्ववित् ।\\
पश्यञ्शृण्वन्स्पृशञ्जिघ्रन्नश्नन्गच्छन्स्वपञ्श्वसन् ॥ ५.८ ॥}
\translate{naiva kiñcitkarōmīti yuktō manyēta tattvavit ।\\
paśyañśr̥ṇvanspr̥śañjighrannaśnangacchansvapañśvasan ॥ 5.8 ॥}
\cquote{A person who is united with Karma Yoga, knows the truth and believes that they do nothing at all. Even though they are engaged in various daily activities such as seeing, hearing, touching, smelling, eating, walking, sleeping, or sitting, they remain detached and are unaffected by these actions.}


\slcol{\Index{प्रलपन् विसृजन्गृह्णन्नु}न्मिषन्निमिषन्नपि ।\\
इन्द्रियाणीन्द्रियार्थेषु वर्तन्त इति धारयन् ॥ ५.९ ॥}
\translate{pralapan visr̥jangr̥hṇannunmiṣannimiṣannapi ।\\
indriyāṇīndriyārthēṣu vartanta iti dhārayan ॥ 5.9 ॥}
\cquote{Even though the senses engage in actions such as speaking, letting go, holding, or opening and closing the eyes, the person who is grounded in truth understands that these activities should be observed and does not get attached to them.}

\newpage
\begin{mananam}{\mananamfont {Mananam sloka - 8, 9}} 
\mananamtext How can I learn to function in this world without being identified with various actions? Can I create a space between the operations of my sense organs and my detached awareness of them? Or, can I observe that there is a sense of ‘I am’, followed by the subsequent arising sense of ‘I am the doer’?\\ 
Can I cultivate the lived knowledge that what I see is a neutral object until I become the seer who has a positive or negative mental reaction to what I see? Can I observe this in my daily life while noticing how each of these (mere seeing versus prejudiced seeing) leave a different impression on the mind? 
\end{mananam}
\WritingHand\enspace{\selfReflect\small{Self Reflection}}
\begin{inspiration}{\mananamfont Inspiration}
\mananamtext We can all benefit by learning the art of holding onto our Self-awareness as we go about our daily activities. It is a means to do all our work and face all our experiences without getting unduly excited or stressed. There is ever a choice within us as to how much we get personally identified and invested in our activities. And even if we do get identified with all of our works, how fast can we withdraw back to our natural peaceful state of Beingness?
\end{inspiration}
\newpage


\slcol{\Index{ब्रह्मण्याधाय कर्माणि} सङ्गं त्यक्त्वा करोति यः ।\\
लिप्यते न स पापेन पद्मपत्रमिवाम्भसा ॥ ५.१० ॥}
\translate{brahmaṇyādhāya karmāṇi saṅgaṁ tyaktvā karōti yaḥ ।\\
lipyatē na sa pāpēna padmapatramivāmbhasā ॥ 5.10 ॥}
\cquote{A person who performs actions by dedicating them to Brahman and relinquishing all attachments is not affected by sin, just like a lotus petal that is not touched by water.}

\slcol{\Index{कायेन मनसा बुद्ध्या} केवलैरिन्द्रियैरपि ।\\
योगिनः कर्म कुर्वन्ति सङ्गं त्यक्त्वात्मशुद्धये ॥ ५.११ ॥}
\translate{kāyēna manasā buddhyā kēvalairindriyairapi ।\\
yōginaḥ karma kurvanti saṅgaṁ tyaktvātmaśuddhayē ॥ 5.11 ॥}
\cquote{Even though Yogis engage in actions using their body, mind, intellect, and senses, they perform these actions selflessly and without any attachment. The actions are performed solely for the purpose of self-purification (ātmaśuddhayē).}

\slcol{\Index{युक्तः कर्मफलं त्यक्त्वा} शान्तिमाप्नोति नैष्ठिकीम् ।\\
अयुक्तः कामकारेण फले सक्तो निबध्यते ॥ ५.१२ ॥}
\translate{yuktaḥ karmaphalaṁ tyaktvā śāntimāpnōti naiṣṭhikīm ।\\
ayuktaḥ kāmakārēṇa phalē saktō nibadhyatē ॥ 5.12 ॥}
\cquote{A disciplined sadhaka (practitioner) who renounces the fruits of their actions attains ultimate peace. However, the undisciplined person, driven by selfish desires and attachment to the results of their actions, is bound to the cycle of Karma.}


\slcol{\Index{सर्वकर्माणि मनसा} संन्यस्यास्ते सुखं वशी ।\\
नवद्वारे पुरे देही नैव कुर्वन्न कारयन् ॥ ५.१३ ॥}
\translate{sarvakarmāṇi manasā saṁnyasyāstē sukhaṁ vaśī ।\\
navadvārē purē dēhī naiva kurvanna kārayan ॥ 5.13 ॥}
\cquote{A self-controlled and disciplined person, having mentally relinquished all actions, peacefully resides within the city of nine gates (the body), neither performing actions nor prompting others (body and senses) to act.}

\newpage
\begin{mananam}{\mananamfont {Mananam sloka - 11}} 
\mananamtext What is the nature of the ego in actions? How does it manifest in my life? How can I attain purification of the ego? Is it necessary to renounce all actions that trigger the ego or is there a balanced approach? What does it mean that yogis are able to function with detachment? Are they completely free from ego or could there be some kind of a functional ego? What is the nature of the latter? 
\end{mananam}
\WritingHand\enspace{\selfReflect\small{Self Reflection}}
\begin{inspiration}{\mananamfont Inspiration}
\mananamtext All the ancient traditions of India emphasized the need for selfless actions as a means of purification of the mind. Seva and karma yoga are invaluable for a sadhak (spiritual practitioner), regardless of what kind of sadhana one embraces. Without it, it is hard to get rid of the ego that functions at different levels and manifests in different ways in a sadhak’s life.
\end{inspiration}
\newpage



\newpage
\begin{mananam}{\mananamfont {Mananam sloka - 13}} 
\mananamtext Am I aware of how the gates of body, such as eyes and ears, are the medium for our interaction with the outside world? Am I aware that what comes in through these gates, and what goes out, as my verbal and physical responses, create an impression on the mind? Am I aware of the moment of choice before my attention and energy goes out through these gates and gets completely identified with the objects of attention?
\end{mananam}
\WritingHand\enspace{\selfReflect\small{Self Reflection}}
\begin{inspiration}{\mananamfont Inspiration}
\mananamtext The inner being within us is the higher Self. Its interactions with the outside world causes us to lose our true identity, leading to various challenges in our interactions with others. To counter this, we must cultivate Self-awareness or the witnessing consciousness as we engage in daily activities. Start with small steps: observe a flower or a bird in the sky, listen to distant sounds, walk within your home or on your terrace, and pay attention to your own breath. In all these moments, practice maintaining partial awareness of your inner being or the higher Self.
\end{inspiration}
\newpage


\slcol{\Index{न कर्तृत्वं न कर्माणि} लोकस्य सृजति प्रभुः ।\\
न कर्मफलसंयोगं स्वभावस्तु प्रवर्तते ॥ ५.१४ ॥}
\translate{na kartr̥tvaṁ na karmāṇi lōkasya sr̥jati prabhuḥ ।\\
na karmaphalasaṁyōgaṁ svabhāvastu pravartatē ॥ 5.14 ॥}
\cquote{The Supreme Being (the Lord) neither creates the sense of doership nor any actions in the world. The connection between actions and their result arises naturally from the inherent nature of individuals and unfolds on its own.}


\slcol{\Index{नादत्ते कस्यचित्पापं} न चैव सुकृतं विभुः ।\\
अज्ञानेनावृतं ज्ञानं तेन मुह्यन्ति जन्तवः ॥ ५.१५ ॥} 
\translate{nādattē kasyacitpāpaṁ na caiva sukr̥taṁ vibhuḥ ।\\
ajñānēnāvr̥taṁ jñānaṁ tēna muhyanti jantavaḥ ॥ 5.15 ॥}
\cquote{The Supreme Being assigns neither sin (demerit) nor virtue (merit) to anyone. It is the veil of ignorance over the true knowledge that causes the living beings to become confused and deluded.}

\slcol{\Index{ज्ञानेन तु तदज्ञानं} येषां नाशितमात्मनः ।\\
तेषामादित्यवज्ज्ञानं प्रकाशयति तत्परम् ॥ ५.१६ ॥}
\translate{jñānēna tu tadajñānaṁ yēṣāṁ nāśitamātmanaḥ ।\\
tēṣāmādityavajjñānaṁ prakāśayati tatparam ॥ 5.16 ॥}
\cquote{Through knowledge, the veil of ignorance that covers the soul is removed. For those who have attained this knowledge, the supreme truth shines forth, illuminating and revealing everything like the sun.}

\slcol{\Index{तद्बुद्धयस्तदात्मान}स्तन्निष्ठास्तत्परायणाः ।\\
गच्छन्त्यपुनरावृत्तिं ज्ञाननिर्धूतकल्मषाः ॥५.१७ ॥} 
\translate{tadbuddhayastadātmānastanniṣṭhāstatparāyaṇāḥ ।\\
gacchantyapunarāvr̥ttiṁ jñānanirdhūtakalmaṣāḥ ॥5.17 ॥}
\cquote{Those whose intellect (buddhi), soul (ātman), and dedication (niṣṭhā) are firmly fixed in that supreme knowledge or truth, and who are completely devoted to it, attain liberation and transcend the cycle of birth and death. Through knowledge, their sins and impurities are completely purified.}

\newpage
\begin{mananam}{\mananamfont {Mananam sloka - 16, 17}} 
\mananamtext Why do the ancient Scriptures tell us we are ignorant? What areas of knowledge elude me despite my academic education and civilised upbringing? Am I capable of discerning the distinction between worldly knowledge and spiritual wisdom? Which knowledge can lead me to the Supreme Truth? What practices can help align my mind with that Ultimate Reality? In what ways can these practices transform my daily life? 
\end{mananam}
\WritingHand\enspace{\selfReflect\small{Self Reflection}}
\begin{inspiration}{\mananamfont Inspiration}
\mananamtext The culmination of all spiritual paths is the eradication of ignorance. Just as darkness yields to light, knowledge alone can dispel ignorance. One insidious effect of ignorance is its concealment of our true identity. Consequently, all our actions arise from this darkness. Continued efforts to establish our intellect in the highest wisdom breaks down all limitations.
\end{inspiration}
\newpage


\slcol{\Index{विद्याविनयसम्पन्ने} ब्राह्मणे गवि हस्तिनि ।\\
शुनि चैव श्वपाके च पण्डिताः समदर्शिनः ॥ ५.१८ ॥}
\translate{vidyāvinayasampannē brāhmaṇē gavi hastini ।\\
śuni caiva śvapākē ca paṇḍitāḥ samadarśinaḥ ॥ 5.18 ॥}
\cquote{The wise, blessed with both knowledge and humility, perceive all beings equally - whether they are a Brahmana (scholar), a cow, an elephant, a dog or even a low-born person (such as a dog-eater). They regard everyone with an equal eye, impartially and without discrimination.}

\slcol{\Index{इहैव तैर्जितः सर्गो} येषां साम्ये स्थितं मनः ।\\
निर्दोषं हि समं ब्रह्म तस्माद्ब्रह्मणि ते स्थिताः ॥ ५.१९ ॥} 
\translate{ihaiva tairjitaḥ sargō yēṣāṁ sāmyē sthitaṁ manaḥ ।\\
nirdōṣaṁ hi samaṁ brahma tasmādbrahmaṇi tē sthitāḥ ॥ 5.19 ॥}
\cquote{In this world and within this very life, those whose minds remain unwaveringly balanced in equanimity have transcended the cycle of birth and death. They are established in Brahman, who is the Supreme, flawless and impartial to all.}


\slcol{\Index{न प्रहृष्येत्प्रियं} प्राप्य नोद्विजेत्प्राप्य चाप्रियम् ।\\
स्थिरबुद्धिरसंमूढो ब्रह्मविद्ब्रह्मणि स्थितः ॥ ५.२० ॥}
\translate{na prahr̥ṣyētpriyaṁ prāpya nōdvijētprāpya cāpriyam ।\\
sthirabuddhirasaṁmūḍhō brahmavidbrahmaṇi sthitaḥ ॥ 5.20 ॥}
\cquote{One who truly knows Brahman remains steadfast and firmly established in Him. Such a person neither rejoices in happiness nor becomes disturbed by unpleasantness. Their intellect remains unwavering and free from delusion.}

\slcol{\Index{बाह्यस्पर्शेष्वसक्तात्मा} विन्दत्यात्मनि यत्सुखम् ।\\
स ब्रह्मयोगयुक्तात्मा सुखमक्षयमश्नुते ॥ ५.२१ ॥}
\translate{bāhyasparśēṣvasaktātmā vindatyātmani yatsukham ।\\
sa brahmayōgayuktātmā sukhamakṣayamaśnutē ॥ 5.21 ॥}
\cquote{A person connected to Brahman through Yoga (meditation) attains eternal and imperishable joy. They are detached from  external sense objects, and find bliss within their own Self-experiences.}


\newpage
\begin{mananam}{\mananamfont {Mananam sloka - 18}} 
\mananamtext By what criteria do I evaluate people in society – based on their status, riches, education, beauty, or achievements? Whom do I honour more and whom do I honour less? Do I hold some people in lower regard? How can I benefit by giving my attention and care to those who are needy? Who are the role models that I can look up to, those whose values and life experiences may elevate my intellectual and spiritual growth?
\end{mananam}
\WritingHand\enspace{\selfReflect\small{Self Reflection}}
\begin{inspiration}{\mananamfont Inspiration}
\mananamtext An exalted human being is one whose response remains equanimous toward all, neither awed by the attention of a rich, powerful individual nor disdainful of someone with low social status.\\
Equanimity should not be mistaken for indifference. Instead, we must learn to recognize those who require our attention and those who deserve reverence. As a practice in shaping our value system, we must shift our focus from our impulse of honouring the wealthy and powerful to prioritising our respect for those who are spiritually wise. Such a value system is a sign of a spiritually advanced society.
\end{inspiration}
\newpage



\newpage
\begin{mananam}{\mananamfont {Mananam sloka - 20}} 
\mananamtext When faced with pleasant experiences, how do I respond? How do I handle unpleasant situations? Am I convinced that rising above the seeking of pleasure and the avoidance of pain is beneficial for my mental well-being? If so, how?\\
Am I conscious of the choices between immediate pleasure and the long-term good in all aspects of life? What is the value of the calm acceptance of results or outcomes, especially after I have put in significant efforts?
\end{mananam}
\WritingHand\enspace{\selfReflect\small{Self Reflection}}
\begin{inspiration}{\mananamfont Inspiration}
\mananamtext What is often seen as a ‘natural impulse’ for humans is regarded as harmful by wise sages and yogis. Almost everyone rushes to seek what is pleasant and avoid what is unpleasant. However, it is precisely these likes and dislikes that lead to inner instability. Cultivating mental equanimity by accepting both pleasant and unpleasant experiences is a highly beneficial practice.
\end{inspiration}
\newpage


\slcol{\Index{ये हि संस्पर्शजा भोगा} दुःखयोनय एव ते ।\\
आद्यन्तवन्तः कौन्तेय न तेषु रमते बुधः ॥ ५.२२ ॥}
\translate{yē hi saṁsparśajā bhōgā duḥkhayōnaya ēva tē ।\\
ādyantavantaḥ kauntēya na tēṣu ramatē budhaḥ ॥ 5.22 ॥}
\cquote{O Kunteya* (Arjuna, son of Kunti), the one who seeks pleasure in sensory experiences that arise from contact is bound to sorrow, as these pleasures are temporary and have both a beginning and an end. In contrast, a wise person does not seek refuge in pleasure for happiness.}

\slcol{\Index{शक्नोतीहैव यः सोढुं} प्राक्छरीरविमोक्षणात् ।\\
कामक्रोधोद्भवं वेगं स युक्तः स सुखी नरः ॥ ५.२३ ॥}
\translate{śaknōtīhaiva yaḥ sōḍhuṁ prākcharīravimōkṣaṇāt ।\\
kāmakrōdhōdbhavaṁ vēgaṁ sa yuktaḥ sa sukhī naraḥ ॥ 5.23 ॥}
\cquote{A person who can withstand the impulses of desire and anger in this life, even before attaining liberation from the body, is a disciplined Yogi. Such a person is truly happy and content.}


\slcol{\Index{योऽन्तःसुखोऽन्तराराम}स्तथान्तर्ज्योतिरेव यः ।\\
स योगी ब्रह्मनिर्वाणं ब्रह्मभूतोऽधिगच्छति ॥ ५.२४ ॥}
\translate{yō:'ntaḥsukhō:'ntarārāmastathāntarjyōtirēva yaḥ ।\\
sa yōgī brahmanirvāṇaṁ brahmabhūtō:'dhigacchati ॥ 5.24 ॥}
\cquote{A person who experiences inner happiness, finds joy within themselves, and whose inner light is spiritual wisdom - such a yogi, is one with Brahman, attains liberation in Brahman.}


\slcol{\Index{लभन्ते ब्रह्मनिर्वाणमृषयः} क्षीणकल्मषाः ।\\ 
छिन्नद्वैधा यतात्मानः सर्वभूतहिते रताः ॥ ५.२५ ॥}
\translate{labhantē brahmanirvāṇamr̥ṣayaḥ kṣīṇakalmaṣāḥ ।\\
chinnadvaidhā yatātmānaḥ sarvabhūtahitē ratāḥ ॥ 5.25 ॥}
\cquote{The sages, whose sins have been eradicated, whose doubts have been dispelled, whose minds are disciplined, and who are dedicated to the well-being of all others, achieve liberation in Brahman.}

\newpage
\begin{mananam}{\mananamfont {Mananam sloka - 21, 22}} 
\mananamtext Am I conscious of the interaction between my senses and sense-objects? For instance, when my eyes perceive something attractive, am I aware of the mental reactions it triggers—such as the desire to acquire or to possess?\\
Do external objects—like food, entertainment, or social interaction—truly bring lasting happiness? Or are they only capable of providing short-term satisfaction? Am I aware of the impermanence inherent in seeking happiness solely from external sources?
\end{mananam}
\WritingHand\enspace{\selfReflect\small{Self Reflection}}
\begin{inspiration}{\mananamfont Inspiration}
\mananamtext We often seek happiness in impermanent external objects. Consequently, this happiness, too, becomes fleeting and leads to suffering when those objects vanish. Our inner thoughts and emotions are closely tied to the external world. Happiness and sorrow are experienced within our minds.\\
Seeing this clearly gives one the power to disconnect the outer world from the inner world. To learn to be happy at will, regardless of outer circumstances, is an essential skill on the spiritual path.
\end{inspiration}
\newpage


\slcol{\Index{कामक्रोधवियुक्तानां} यतीनां यतचेतसाम् ।\\
अभितो ब्रह्मनिर्वाणं वर्तते विदितात्मनाम् ॥ ५.२६ ॥} 
\translate{kāmakrōdhaviyuktānāṁ yatīnāṁ yatacētasām ।\\
abhitō brahmanirvāṇaṁ vartatē viditātmanām ॥ 5.26 ॥}
\cquote{Those who have overcome desire and anger, who have controlled their minds and who practice self-discipline, attain liberation in Brahman in all situations. This state can only be achieved by a self-realized individual.}

\slcol{\Index{स्पर्शान्कृत्वा} बहिर्बाह्यांश्चक्षुश्चैवान्तरे भ्रुवोः ।\\
प्राणापानौ समौ कृत्वा नासाभ्यन्तरचारिणौ ॥ ५.२७ ॥}
\translate{sparśānkr̥tvā bahirbāhyāṁścakṣuścaivāntarē bhruvōḥ ।\\
prāṇāpānau samau kr̥tvā nāsābhyantaracāriṇau ॥ 5.27 ॥}
\cquote{After shutting off the external sensory channels, such as touch, and directing the eyes inward, between the eyebrows, one harmonizes the breath - prāṇa* (inhalation) and apāna* (exhalation) - guiding them through the nostrils. This is the technique of controlling the senses and vital breath during meditation.}


\slcol{\Index{यतेन्द्रियमनोबुद्धि}र्मुनिर्मोक्षपरायणः ।\\ 
विगतेच्छाभयक्रोधो यः सदा मुक्त एव सः ॥ ५.२८ ॥}
\translate{yatēndriyamanōbuddhirmunirmōkṣaparāyaṇaḥ ।\\ 
vigatēcchābhayakrōdhō yaḥ sadā mukta ēva saḥ ॥ 5.28 ॥} 
\cquote{A person who has controlled their senses, mind and intellect, with liberation as their supreme goal, and who has let go of desire, fear, and anger, remains eternally free and liberated.}

\slcol{\Index{भोक्तारं यज्ञतपसां} सर्वलोकमहेश्वरम् ।\\ 
सुहृदं सर्वभूतानां ज्ञात्वा मां शान्तिमृच्छति ॥ ५.२९ ॥ }
\translate{bhōktāraṁ yajñatapasāṁ sarvalōkamahēśvaram ।\\
suhr̥daṁ sarvabhūtānāṁ jñātvā māṁ śāntimr̥cchati ॥ 5.29 ॥} 
\cquote{One who recognizes Me as the ultimate enjoyer of all sacrifices and penances, the supreme ruler of all realms, and the best friend of all living beings, attains lasting peace.}

\newpage
\begin{mananam}{\mananamfont {Mananam sloka - 26}} 
\mananamtext What role do desires and related emotions play in my life? Are my desires healthy and constructive for my long-term well-being? Is my motivation in life driven solely by self-centred desires? Can I find motivation to act based on the intention of the well-being of others? Can I recall instances where strong attachment led to anger due to perceived obstruction of my desires? What does it mean to me to attain the state of the sages, devoid of all passions and cravings? 
\end{mananam}
\WritingHand\enspace{\selfReflect\small{Self Reflection}}
\begin{inspiration}{\mananamfont Inspiration}
\mananamtext The common person’s concept of freedom involves having all his needs and desires fulfilled as soon as they arise in his mind. In contrast, yogis attain freedom through vigilant mental awareness, where they allow no unwanted desires for sensory gratification or egoic validation. If any desires do arise, they are immediately transmuted and directed toward the well-being of others. The very act of desiring ties us down with our body identification. True freedom is to abide in the soul’s natural state of supreme bliss.
\end{inspiration}
\newpage

\begin{center}
  ॐ तत्सदिति श्रीमद्भगवद्गीतासूपनिषत्सु ब्रह्मविद्यायां योगशास्त्रे\\
  श्रीकृष्णार्जुनसम्वादे संन्यासयोगो नाम पञ्चमोऽध्यायः॥\\
  \arialUC\small ōṃ tatsaditi śrīmadbhagavadgītās ūpaniṣatsu \\brahmavidyāyāṃ yōgaśāstrē\\ śrīkṛṣṇārjunasaṃvādē sannyāsa-yoga\\nāma sannyāsa-yoga॥
\end{center}
\chapEnd